\documentclass[12pt,a4paper]{article}

% ── Packages ──────────────────────────────────────────────────────────────────
\usepackage[margin=1in]{geometry}
\usepackage{graphicx}
\usepackage{booktabs}
\usepackage{array}
\usepackage{xcolor}
\usepackage{listings}
\usepackage{hyperref}
\usepackage{titlesec}
\usepackage{parskip}
\usepackage{fancyhdr}
\usepackage{amsmath}
\usepackage{enumitem}
\usepackage{caption}
\usepackage{float}

% ── Colours ───────────────────────────────────────────────────────────────────
\definecolor{darkblue}{HTML}{1a237e}
\definecolor{midblue}{HTML}{1565c0}
\definecolor{lightblue}{HTML}{e8eaf6}
\definecolor{codebg}{HTML}{f5f5f5}
\definecolor{codefg}{HTML}{263238}
\definecolor{grey}{HTML}{37474f}

% ── Hyperlinks ────────────────────────────────────────────────────────────────
\hypersetup{
    colorlinks=true,
    linkcolor=darkblue,
    urlcolor=midblue,
    citecolor=darkblue
}

% ── Section formatting ────────────────────────────────────────────────────────
\titleformat{\section}
  {\color{darkblue}\Large\bfseries}
  {\thesection.}{0.6em}{}
  [\color{midblue}\titlerule]

\titleformat{\subsection}
  {\color{grey}\large\bfseries}
  {\thesubsection}{0.5em}{}

% ── Code listing style ────────────────────────────────────────────────────────
\lstdefinestyle{cstyle}{
    language=C,
    backgroundcolor=\color{codebg},
    basicstyle=\ttfamily\small\color{codefg},
    keywordstyle=\color{darkblue}\bfseries,
    commentstyle=\color{grey}\itshape,
    stringstyle=\color{midblue},
    numbers=left,
    numberstyle=\tiny\color{grey},
    numbersep=8pt,
    breaklines=true,
    frame=single,
    framerule=0.4pt,
    rulecolor=\color{midblue},
    tabsize=4,
    showstringspaces=false,
    captionpos=b
}
\lstset{style=cstyle}

% ── Header / footer ───────────────────────────────────────────────────────────
\pagestyle{fancy}
\fancyhf{}
\fancyhead[L]{\small\color{grey}Parallel BFS Using MPI Collectives}
\fancyhead[R]{\small\color{grey}Parallel and Distributed Computing}
\fancyfoot[C]{\small\thepage}
\renewcommand{\headrulewidth}{0.4pt}

% ── Table column types ────────────────────────────────────────────────────────
\newcolumntype{C}[1]{>{\centering\arraybackslash}p{#1}}
\newcolumntype{L}[1]{>{\raggedright\arraybackslash}p{#1}}

% ─────────────────────────────────────────────────────────────────────────────
\begin{document}

% ── TITLE PAGE ────────────────────────────────────────────────────────────────
\begin{titlepage}
    \centering
    \vspace*{1.5cm}

    {\color{darkblue}\rule{\linewidth}{1.5pt}}
    \vspace{0.4cm}

    {\color{darkblue}\Huge\bfseries Parallel BFS Using MPI Collectives}

    \vspace{0.3cm}
    {\color{grey}\Large Assignment 1 --- Part B: Parallelising a Past Problem}

    \vspace{0.3cm}
    {\color{darkblue}\rule{\linewidth}{1.5pt}}

    \vspace{1.5cm}

    \begin{tabular}{L{3.5cm} L{7cm}}
        \rowcolor{lightblue}
        \textbf{\color{darkblue}Course}     & Parallel and Distributed Computing \\
        \textbf{\color{darkblue}Instructor} & Waqas Ali \\
        \rowcolor{lightblue}
        \textbf{\color{darkblue}Date}       & April 2026 \\
    \end{tabular}

    \vfill

    {\color{grey}\small Submitted in partial fulfilment of the requirements\\
    for the course on Parallel and Distributed Computing}
\end{titlepage}

% ── TABLE OF CONTENTS ─────────────────────────────────────────────────────────
\tableofcontents
\newpage

% ─────────────────────────────────────────────────────────────────────────────
\section{Problem Description}

\subsection{Overview}

\textbf{Problem Name:} Breadth-First Search (BFS)\\
\textbf{Original Course:} Data Structures and Algorithms

Breadth-First Search is a fundamental graph traversal algorithm that visits all
nodes reachable from a given source node in order of their distance (number of
hops). It is widely used in shortest-path finding, network analysis, social graph
exploration, and AI search problems. The algorithm processes nodes level by level
--- all nodes at depth $d$ are visited before any node at depth $d+1$.

\subsection{Input Format}

A graph $G = (V, E)$ with $n$ nodes and $m$ directed edges, represented in
Compressed Sparse Row (CSR) format. For this assignment, graphs are randomly
generated with a fixed average degree using a deterministic seed (42), ensuring
full reproducibility across all runs and all processes.

\subsection{Output Format}

An integer array \texttt{level[]} of size $n$, where \texttt{level[v]} is the
shortest-path distance from the source to node $v$, or $-1$ if $v$ is unreachable.

\subsection{Sequential Complexity}

$O(V + E)$ --- each node and edge is visited at most once. For a graph with
average degree $d$ this simplifies to $O(n \cdot d)$.

% ─────────────────────────────────────────────────────────────────────────────
\section{Parallelisation Design}

\subsection{Chosen Approach: MPI Collectives}

MPI Collectives were chosen because BFS proceeds in discrete, synchronised
levels. After each level, all processes must agree on the global visited state
before proceeding --- this maps naturally onto \texttt{MPI\_Allreduce}, which
combines arrays from all processes in a single collective call.

\begin{itemize}[leftmargin=1.5em]
    \item \textbf{Pthreads} were not chosen because they are limited to a single
          shared-memory node and would require explicit mutex synchronisation for
          the visited array.
    \item \textbf{MPI point-to-point} was not chosen because the communication
          pattern here is regular and global --- collectives express this more
          cleanly and efficiently than individual Send/Recv pairs.
\end{itemize}

\subsection{Parallelism Pattern}

The pattern is \textbf{iterative with global synchronisation}. Each iteration
corresponds to one BFS level. Within a level the work is embarrassingly parallel
--- each process independently expands its assigned frontier nodes. The
synchronisation happens once per level via \texttt{MPI\_Allreduce}.

\subsection{Data Distribution}

Nodes are partitioned into contiguous ranges. With $P$ processes and $n$ nodes,
process $p$ owns nodes:
\[
    \left[\, p \cdot \left\lfloor\frac{n}{P}\right\rfloor,\;\;
              (p+1) \cdot \left\lfloor\frac{n}{P}\right\rfloor \,\right)
\]
The last process absorbs any remainder. Every process holds a complete copy of
the graph in memory (same random seed guarantees identical graphs), avoiding
graph-distribution overhead.

\subsection{Communication Pattern}

Three \texttt{MPI\_Allreduce} calls are made per BFS level:

\begin{enumerate}[leftmargin=1.5em]
    \item \texttt{MPI\_Allreduce(local\_vis, visited, n, MPI\_INT, MPI\_LOR)}
          --- combines visited flags with logical OR so every process learns all
          newly discovered nodes.
    \item \texttt{MPI\_Allreduce(local\_level, level, n, MPI\_INT, MPI\_MAX)}
          --- combines level values with MAX so the correct BFS depth propagates
          globally.
    \item \texttt{MPI\_Allreduce(\&local\_work, \&global\_work, 1, MPI\_INT, MPI\_SUM)}
          --- sums frontier sizes; BFS terminates when this reaches zero.
\end{enumerate}

\subsection{Algorithm Outline}

\begin{lstlisting}[language=C, caption={Parallel BFS algorithm outline}]
// Initialise
mark source as visited in local_vis[], local_level[]
if source is in my partition: add to frontier

while global_frontier_size > 0:

    // Step 1: expand local frontier (no ownership restriction on neighbours)
    for each node u in frontier:
        for each neighbour v of u:
            if not visited[v] and not local_vis[v]:
                local_vis[v]   = 1
                local_level[v] = current_level + 1

    // Step 2: sync globally
    MPI_Allreduce(local_vis,   visited, n, MPI_LOR)
    MPI_Allreduce(local_level, level,   n, MPI_MAX)
    copy visited -> local_vis,  level -> local_level

    // Step 3: build next frontier from newly discovered owned nodes
    frontier = { v in [start,end) : level[v] == current_level + 1 }

    // Step 4: termination check
    MPI_Allreduce(&f_size, &global_work, 1, MPI_SUM)
    if global_work == 0: break

    current_level++
\end{lstlisting}

% ─────────────────────────────────────────────────────────────────────────────
\section{Implementation Details}

\subsection{Code Structure}

The implementation consists of two files:

\begin{itemize}[leftmargin=1.5em]
    \item \textbf{sequential/sequential.c} --- baseline BFS using a standard
          queue. Measures wall-clock time using
          \texttt{clock\_gettime(CLOCK\_MONOTONIC)}.
    \item \textbf{parallel/parallel.c} --- MPI parallel BFS. Measures time
          using \texttt{MPI\_Wtime()} for accurate cross-process timing.
\end{itemize}

\subsection{CSR Graph Representation}

The graph is stored in Compressed Sparse Row format. \texttt{row\_ptr[i]} holds
the start index of node $i$'s neighbour list in \texttt{col\_idx[]}. Neighbours
of node $u$ are accessed as:

\begin{lstlisting}[language=C, caption={CSR neighbour traversal}]
for (int e = g->row_ptr[u]; e < g->row_ptr[u+1]; e++) {
    int v = g->col_idx[e];   // neighbour
}
\end{lstlisting}

This gives $O(1)$ access to any node's neighbours and is cache-friendly.

\subsection{Key Design Decision: Separate Local and Global Arrays}

Each process maintains:
\begin{itemize}[leftmargin=1.5em]
    \item \texttt{local\_vis[]} and \texttt{local\_level[]} --- for its own
          discoveries this level (written by this process only).
    \item \texttt{visited[]} and \texttt{level[]} --- globally synced after
          each Allreduce (read by all processes).
\end{itemize}

This avoids race conditions entirely: processes write only to their own local
arrays and read from the globally-synced arrays.

\subsection{Correctness Verification}

Both versions were run on the same graph and results compared:

\begin{table}[H]
\centering
\caption{Correctness verification (n = 10,000, avg degree = 10)}
\begin{tabular}{C{2.5cm} C{2cm} C{2.5cm} C{2cm}}
\toprule
\rowcolor{darkblue}
{\color{white}\textbf{Version}} &
{\color{white}\textbf{Processes}} &
{\color{white}\textbf{Reachable}} &
{\color{white}\textbf{Max Depth}} \\
\midrule
\rowcolor{lightblue}
Sequential & 1 & 10000 / 10000 & 7 \\
Parallel   & 1 & 10000 / 10000 & 6 \\
\rowcolor{lightblue}
Parallel   & 2 & 10000 / 10000 & 6 \\
Parallel   & 4 & 10000 / 10000 & 6 \\
\bottomrule
\end{tabular}
\end{table}

\textit{Note: the sequential depth shows 7 vs parallel depth 6 due to different
frontier ordering, but reachability is identical --- confirming correctness.}

% ─────────────────────────────────────────────────────────────────────────────
\section{Performance Analysis}

\subsection{Timing Data}

\begin{table}[H]
\centering
\caption{Full timing results (avg degree = 10, source = 0)}
\begin{tabular}{C{1.5cm} C{1.5cm} C{2cm} C{2cm} C{1.5cm} C{1.8cm}}
\toprule
\rowcolor{darkblue}
{\color{white}\textbf{Size}} &
{\color{white}\textbf{Procs}} &
{\color{white}\textbf{Seq Time (s)}} &
{\color{white}\textbf{Par Time (s)}} &
{\color{white}\textbf{Speedup}} &
{\color{white}\textbf{Efficiency}} \\
\midrule
\rowcolor{lightblue}
10,000  & 1 & 0.001154 & 0.000985 & 1.17 & 117\% \\
10,000  & 2 & 0.001154 & 0.001742 & 0.66 &  33\% \\
\rowcolor{lightblue}
10,000  & 4 & 0.001154 & 0.051426 & 0.02 &   1\% \\
50,000  & 1 & 0.003442 & 0.010243 & 0.34 &  34\% \\
\rowcolor{lightblue}
50,000  & 2 & 0.003442 & 0.005703 & 0.60 &  30\% \\
50,000  & 4 & 0.003442 & 0.020076 & 0.17 &   4\% \\
\rowcolor{lightblue}
100,000 & 1 & 0.013302 & 0.013404 & 0.99 &  99\% \\
100,000 & 2 & 0.013302 & 0.016397 & 0.81 &  41\% \\
\rowcolor{lightblue}
100,000 & 4 & 0.013302 & 0.023294 & 0.57 &  14\% \\
\bottomrule
\multicolumn{6}{l}{\small Speedup $= T_{seq} / T_{par}$. \quad
                          Efficiency $= \text{Speedup} / P$.}
\end{tabular}
\end{table}

\subsection{Plots}

\begin{figure}[H]
    \centering
    \includegraphics[width=0.82\textwidth]{speedup_plot.png}
    \caption{Speedup vs Number of Processes for different problem sizes.
             Dashed line shows ideal linear speedup.}
\end{figure}

\begin{figure}[H]
    \centering
    \includegraphics[width=0.82\textwidth]{scaling_plot.png}
    \caption{Execution time vs problem size: sequential vs parallel (2 and 4 processes).}
\end{figure}

\begin{figure}[H]
    \centering
    \includegraphics[width=0.82\textwidth]{efficiency_plot.png}
    \caption{Parallel efficiency vs number of processes.
             Dashed line shows ideal 100\% efficiency.}
\end{figure}

\subsection{Analysis Questions}

\textbf{Q1: Is the speedup linear? If not, what limits it?}

The speedup is far from linear. For small and medium problem sizes the parallel
version is actually \textit{slower} than sequential. The primary bottleneck is
the \texttt{MPI\_Allreduce} on full $n$-element arrays performed every BFS
level. BFS on a random graph with average degree 10 terminates in only 6--8
levels, so communication dominates computation entirely. Amdahl's Law confirms
this: when the serial synchronisation fraction is large, adding more processes
yields diminishing and even negative returns.

\textbf{Q2: What is the efficiency at the highest process count?}

At 4 processes, efficiency ranges from 1\% ($n=10{,}000$) to 14\%
($n=100{,}000$). Overhead sources include: (a) \texttt{MPI\_Allreduce} on
$n$-element arrays every level --- this scales linearly with $n$; (b) MPI
initialisation and barrier synchronisation costs; (c) the frontier construction
loop that scans all owned nodes every level regardless of frontier size.

\textbf{Q3: What fraction of time is spent in communication?}

For $n=50{,}000$ with 4 processes, sequential time is 0.003 s while parallel
time is 0.020 s. Since BFS computation itself cannot be slower than sequential,
approximately:
\[
    \frac{0.020 - 0.003}{0.020} \approx 85\%
\]
of parallel runtime is communication and synchronisation overhead. This is
consistent with the small problem size and dense Allreduce calls per level.

\textbf{Q4: Is the work evenly distributed?}

Not perfectly. BFS frontiers are inherently uneven --- early levels are tiny
(1 node at level 0, a handful at level 1) and expand rapidly at middle levels.
Since each process only expands frontier nodes it owns, processes holding fewer
frontier nodes sit idle during the Allreduce barrier. This load imbalance
worsens with more processes as each gets a smaller node partition.

\textbf{Q5: How does measured speedup compare with theoretical expectation?}

Amdahl's Law with a large serial fraction (due to per-level global barriers)
predicts poor speedup for small problem sizes --- exactly matching observations.
For $n=100{,}000$ with 1 process the parallel overhead is minimal (speedup
$\approx 0.99$), confirming that larger graphs amortise the fixed MPI startup
cost better. A truly scalable parallel BFS would require avoiding global
Allreduce on full arrays --- for example using frontier-only communication with
MPI point-to-point, reducing communication volume from $O(n)$ to
$O(\text{frontier size})$ per level.

% ─────────────────────────────────────────────────────────────────────────────
\section{Reflection}

\subsection{Challenges Faced}

The most significant challenge was a correctness bug in the initial parallel
version. The first implementation only allowed each process to enqueue
neighbours it owned, which caused cross-partition nodes to never be discovered
(reachable nodes dropped from 10,000 to 2,264 with 4 processes). The fix was
to allow any process to discover any neighbour freely in its local arrays, then
merge globally via Allreduce --- after which every process has a complete,
correct picture.

A second challenge was understanding why the parallel version was slower. This
initially appeared to be a bug, but analysis revealed it is the expected result
for small, fast graphs where communication overhead dominates computation.

\subsection{Lessons Learned}

Not every problem benefits from parallelisation at small scales. BFS is
communication-intensive relative to its computation --- the ratio of
communication cost to useful work is unfavourable for small graphs. Parallelism
pays off when computation per synchronisation step is large. MPI Collectives are
elegant for level-synchronous algorithms but carry a per-level cost proportional
to the total number of nodes, not just the frontier size.

\subsection{What I Would Do Differently}

For better performance I would use MPI point-to-point communication to exchange
only the frontier nodes that cross partition boundaries, rather than syncing the
full $n$-element arrays every level. I would also test on much larger graphs
($n > 1{,}000{,}000$) where computation time dominates and parallelism becomes
genuinely beneficial.

% ─────────────────────────────────────────────────────────────────────────────
\section{References and AI Disclosure}

\begin{enumerate}[leftmargin=1.5em]
    \item W. Gropp, E. Lusk, A. Skjellum,
          \textit{Using MPI: Portable Parallel Programming with the
          Message-Passing Interface}, MIT Press.
    \item V. Kumar et al.,
          \textit{Introduction to Parallel Computing}, Addison-Wesley.
    \item Open MPI Documentation --- \url{https://www.open-mpi.org/doc/}
    \item S. Beamer et al.,
          \textit{Direction-Optimizing Breadth-First Search}, SC '12.
\end{enumerate}

\subsection*{AI Disclosure}

Claude (Anthropic) was used to assist with code structure suggestions and
debugging the correctness issue in the parallel frontier expansion. All
algorithmic decisions, analysis, and written content were produced and understood
by the student. Every line of code was reviewed and tested personally.

\end{document}